% mazzi: 03
\chapter{L'eccitabilità cellulare}

\textbf{Cellula}: più piccola unità funzionale dell'organismo; ogni tipo cellulare è specializzato
per compiti specifici, ma tutte le cellule condividono alcune caratteristiche di base.

\textbf{Cellule eccitabili}: cellule capaci di reagire a uno \textit{stimolo adeguato} con una
variazione delle proprietà di membrana.
\begin{itemize}
  \item sono essenzialmente le cellule del sistema nervoso (neuroni) e le cellule muscolari;
  \item strutturalmente uguali a tutte le altre cellule: l'eccitabilità è una proprietà in più,
    non una diversa costituzione.
\end{itemize}

% ---------------------------------------------------------------------------
\section{La membrana cellulare}

\textbf{Membrana plasmatica}: separa l'ambiente intracellulare da quello extracellulare (liquido
interstiziale).

\rubrica{Struttura}
\begin{itemize}
  \item \textbf{bistrato lipidico}: doppio strato di fosfolipidi (modello a \textit{mosaico
    fluido}); isolante, non si fa attraversare da acqua e cariche elettriche $\rightarrow$
    \textit{idrofobico};
  \item \textbf{proteine immerse} nel bistrato:
    \begin{itemize}
      \item \textit{strutturali}: danno la forma alla cellula;
      \item \textit{proteine-canale}: attraversano la membrana da parte a parte e mettono in
        comunicazione ambiente intracellulare ed extracellulare; essendo proteine permettono il
        passaggio di acqua e cariche elettriche $\rightarrow$ \textit{idrofiliche}.
    \end{itemize}
\end{itemize}

\rubrica{Semipermeabilità}
Si fa attraversare da alcune specie ioniche e non da altre; i passaggi avvengono secondo le leggi
chimico-fisiche della \textbf{diffusione}, del \textbf{trasporto facilitato} e del \textbf{trasporto
attivo}.

\rubrica{Canali ionici}
Proteine formate da più subunità proteiche che attraversano il bistrato. Determinano la
caratteristica principale dei due ambienti: una \textbf{diversa distribuzione delle specie
ioniche}.
\begin{itemize}
  \item \textbf{ambiente intracellulare}: ricco di $K^+$ e anioni proteici ($A^-$);
  \item \textbf{ambiente extracellulare}: ricco di $Na^+$, $Ca^{2+}$, $Cl^-$, $HCO_3^-$.
\end{itemize}

% ---------------------------------------------------------------------------
\section{Il potenziale di riposo}

Essendo semipermeabile, la membrana attribuisce a ogni specie ionica un proprio \textbf{potenziale
di equilibrio}: il valore di potenziale al quale quello ione non ha più tendenza netta a spostarsi.

\begin{table}[htbp]
  \centering
  \small
  \renewcommand{\arraystretch}{1.3}
  \begin{tabularx}{\textwidth}{|l|c|c|c|}
    \hline
    \textbf{Ione} & \textbf{Conc.\ intracellulare (mM)} & \textbf{Conc.\ extracellulare (mM)}
      & \textbf{Potenziale di equilibrio (mV)} \\
    \hline
    $K^+$  & 140 & 5   & $-75$ \\
    \hline
    $Na^+$ & 5   & 140 & $+55$ \\
    \hline
    $Cl^-$ & 4,2 & 110 & $-60$ \\
    \hline
    $A^-$  & \multicolumn{3}{c|}{non possono attraversare la membrana} \\
    \hline
  \end{tabularx}
  \caption{Distribuzione delle cariche elettriche ai due lati della membrana.}
  \label{tab:distribuzione-ionica}
\end{table}

\begin{itemize}
  \item il dato più rilevante per il funzionamento cellulare è la differenza di concentrazione di
    $K^+$ (140 dentro contro 5 fuori) e, speculare, di $Na^+$ (5 dentro contro 140 fuori);
  \item i valori della colonna extracellulare sono gli stessi che compaiono nel dosaggio degli
    elettroliti alle analisi del sangue.
\end{itemize}

\rubrica{La membrana come condensatore}
Il bistrato lipidico isolante fa sì che la membrana si comporti da \textbf{condensatore}: separa e
segrega le cariche elettriche sulle due facce.
\begin{itemize}
  \item all'interno si accumulano cariche negative, all'esterno cariche positive;
  \item cariche di segno opposto si attraggono, ma a riposo non possono attraversare la membrana
    perché le proteine-canale sono chiuse;
  \item lontano dalla membrana le cariche $+$ e $-$ si equivalgono: la separazione riguarda solo i
    due sottili strati adiacenti alla superficie interna ed esterna.
\end{itemize}

\textbf{Potenziale di riposo ($V_m$)}: differenza di potenziale che ne risulta, di circa $-70$\,mV
all'interno rispetto all'esterno (alcuni testi riportano $-60$\,mV: conta il concetto, cioè che
l'interno sia nettamente negativo). È la condizione che rende la membrana eccitabile, ossia capace
di rispondere a uno stimolo.

\rubrica{Origine della differenza di potenziale}
Modello sperimentale: due vaschette separate da una membrana selettiva, entrambe con $NaCl$ in
soluzione.
\begin{enumerate}
  \item \textbf{membrana selettiva al $Na^+$}: il $Na^+$ si sposta dalla vaschetta 1 alla 2 secondo
    gradiente, il $Cl^-$ no $\rightarrow$ nella vaschetta 2 prevalgono le cariche positive e si
    genera una differenza di potenziale tra 1 e 2;
  \item \textbf{membrana selettiva al $Cl^-$}: si sposta il $Cl^-$ $\rightarrow$ la vaschetta 2
    accumula cariche negative.
\end{enumerate}
Nella cellula vale lo stesso principio: gli anioni proteici ($A^-$) sono molecole troppo grandi per
attraversare la membrana, anche a canali aperti $\rightarrow$ si accumulano all'interno, che risulta
negativo rispetto all'esterno.

\subsection{Pompa sodio-potassio}

\textbf{Pompa $Na^+/K^+$}: proteina \textit{ATPasica} che attraversa la membrana da parte a parte;
spiega perché $Na^+$ e $K^+$, entrambi monovalenti positivi, si accumulino uno fuori e uno dentro.
\begin{itemize}
  \item \textbf{funzionamento}: espelle 3 $Na^+$ dalla cellula e recupera 2 $K^+$ dall'esterno;
  \item \textbf{ATPasica}: nel suo esercizio consuma ATP (trasporto attivo);
  \item insieme alle proprietà della membrana genera e mantiene la differenza di potenziale di
    $-70$\,mV.
\end{itemize}

\textbf{Glicosidi cardiaci}: farmaci e sostanze tossiche possono interferire con il funzionamento
cellulare \textbf{bloccando la pompa} $\rightarrow$ da qui sia l'efficacia terapeutica in alcune
circostanze sia la tossicità. Argomento ripreso con la fisiologia del cuore.

% ---------------------------------------------------------------------------
\section{Proprietà passive di membrana}

Registrazione dell'attività elettrica di una cellula immersa nel liquido interstiziale: un
\textbf{elettrodo di potenziale} confronta interno ed esterno (lettura a riposo $\sim-70$\,mV),
un \textbf{elettrodo di corrente} collegato a un generatore inietta cariche con uno stimolo
\textit{a onda quadra} di durata e intensità note.

\begin{itemize}
  \item \textbf{depolarizzazione}: iniezione di cariche positive $\rightarrow$ $V_m$ meno negativo
    ($-70 \rightarrow -68 \rightarrow -66$\,mV e così via, in funzione della corrente immessa);
  \item \textbf{iperpolarizzazione}: iniezione di cariche negative $\rightarrow$ $V_m$ più negativo.
\end{itemize}

\textbf{Potenziali elettrotonici}: le variazioni di potenziale così ottenute; sono lente e graduate.
\begin{itemize}
  \item la curva di risposta è \textbf{arrotondata} perché le cariche impiegano un certo tempo a
    entrare nella cellula e a modificarne il potenziale;
  \item le variazioni di voltaggio sono \textbf{proporzionali all'intensità di corrente ($I$)}
    somministrata.
\end{itemize}

\textbf{Modello elettrico}: la membrana si comporta come un circuito dotato di resistenza e
capacità.
\begin{itemize}
  \item \textbf{resistenza}: i canali di membrana, che si fanno attraversare dalle cariche;
  \item \textbf{capacità}: il bistrato lipidico;
  \item in questo regime la membrana segue la \textbf{prima legge di Ohm}: la variazione di
    voltaggio registrata è direttamente proporzionale all'intensità di corrente somministrata, alla
    sua durata e alla resistenza di membrana.
\end{itemize}

% ---------------------------------------------------------------------------
\section{Il potenziale d'azione}

Aumentando progressivamente l'intensità dello stimolo, la membrana a un certo punto \textbf{smette
di seguire la prima legge di Ohm}: si comporta da sistema biologico e non più da circuito passivo.

\textbf{Potenziale di soglia}: valore di $V_m$ intorno a $-55$\,mV oltre il quale scatta la risposta
attiva.
\begin{itemize}
  \item \textbf{stimoli sottosoglia}: generano solo potenziali elettrotonici (risposte passive,
    lente e graduate);
  \item \textbf{stimoli soprasoglia}: generano il \textbf{potenziale d'azione}, risposta
    \textit{tutto o nulla}.
\end{itemize}

\rubrica{Fasi del potenziale d'azione}
\begin{enumerate}
  \item \textbf{prepotenziale}: dipende dallo stimolo elettrico o chimico somministrato; porta
    $V_m$ verso la soglia;
  \item \textbf{fase ascendente, sodio-dipendente} (\textit{depolarizzazione}): raggiunta la soglia
    si aprono le proteine-canale per il $Na^+$ \textit{voltaggio-dipendenti} $\rightarrow$ aumenta
    la conduttanza al sodio ($g_{Na}$) e il $Na^+$ entra secondo gradiente fino al proprio
    potenziale di equilibrio; si supera lo 0 e l'interno diventa positivo (\textbf{inversione});
  \item \textbf{fase discendente, potassio-dipendente} (\textit{ripolarizzazione}): si aprono i
    canali voltaggio-dipendenti per il $K^+$; il $K^+$, più concentrato all'interno, esce secondo
    gradiente $\rightarrow$ uscita di cariche positive e ripolarizzazione fino al potenziale di
    equilibrio del potassio;
  \item \textbf{iperpolarizzazione postuma} e ritorno al riposo: chiusi anche i canali del $K^+$,
    la pompa $Na^+/K^+$ riprende il sopravvento e riporta la membrana al potenziale di riposo,
    pronta per un nuovo potenziale d'azione.
\end{enumerate}
Durata dell'intero evento: nell'ordine del millisecondo.

\rubrica{Il potenziale d'azione come informazione}
Costituisce la risposta della cellula allo stimolo ed è esso stesso un'\textbf{informazione}. Il
sistema nervoso comprende soltanto stimoli elettrici: gli stimoli di altra natura (chimici,
elettromagnetici, pressori) vanno tramutati in stimoli elettrici perché possano essere compresi,
integrati e tradotti in un comando all'organo effettore.

\subsection{Periodi refrattari}

\begin{itemize}
  \item \textbf{periodo refrattario assoluto}: durante il potenziale d'azione non si può generare
    un altro potenziale d'azione, qualunque sia l'intensità dello stimolo;
  \item \textbf{periodo refrattario relativo}: è possibile generare un potenziale d'azione solo
    aumentando notevolmente l'intensità dello stimolo. Corrisponde al tempo necessario alla pompa
    $Na^+/K^+$ per ributtare fuori tutto il $Na^+$ entrato durante la depolarizzazione e rimettere
    dentro tutto il $K^+$ uscito durante la ripolarizzazione.
\end{itemize}

% ---------------------------------------------------------------------------
\section{Caliemia ed eccitabilità}

Per l'eccitabilità cellulare è fondamentale mantenere la distribuzione elettrolitica descritta; la
concentrazione plasmatica di $K^+$ ne è il determinante più diretto.
\begin{itemize}
  \item \textbf{normocaliemia} ($K^+$ plasmatico normale, 5\,mM): uno stimolo sottosoglia genera
    solo un potenziale graduato; uno stimolo soprasoglia scatena il potenziale d'azione;
  \item \textbf{ipercaliemia} (aumento del $K^+$ plasmatico, p.\,es.\ 6--7\,mM): il potenziale di
    riposo si avvicina alla soglia $\rightarrow$ la cellula diventa \textbf{più eccitabile}, e uno
    stimolo che sarebbe normalmente sottosoglia può far partire un potenziale d'azione;
  \item \textbf{ipocaliemia} (diminuzione del $K^+$ plasmatico): la membrana si iperpolarizza
    $\rightarrow$ la cellula diventa \textbf{meno eccitabile} e per raggiungere la soglia occorre
    aumentare l'intensità dello stimolo; si riduce la probabilità che il neurone risponda a uno
    stimolo normalmente soprasoglia.
\end{itemize}

% ---------------------------------------------------------------------------
\section{Propagazione del potenziale lungo l'assone}

Il potenziale d'azione è un'informazione: deve viaggiare lungo l'assone e raggiungere il terminale
assonico per essere trasmesso ad altre cellule. La conduzione dell'impulso può essere di due tipi.

\subsection{Conduzione punto a punto}

Tipica degli \textbf{assoni non mielinizzati}; \textbf{lenta}, perché il potenziale va condotto
punto per punto lungo un assone che può essere lungo anche un metro.
\begin{itemize}
  \item nel punto stimolato si aprono i canali del $Na^+$ $\rightarrow$ \textbf{zona
    depolarizzata}: cariche positive all'interno, negative all'esterno;
  \item si formano \textbf{correnti di circuito locale}: le cariche positive interne viaggiano
    lungo l'assone e depolarizzano il punto successivo (\textit{zona attiva}), mentre all'esterno le
    cariche negative rifluiscono verso il punto che si era depolarizzato;
  \item il punto precedente torna così negativo all'interno e positivo all'esterno (\textit{zona
    refrattaria}, canali del $Na^+$ inattivati) e il potenziale procede lungo tutto l'assone.
\end{itemize}

\subsection{Conduzione saltatoria}

Tipica dei neuroni ricoperti da \textbf{guaina mielinica}; lo scambio di cariche elettriche avviene
solo in punti determinati, i \textbf{nodi di Ranvier}.
\begin{itemize}
  \item un nodo si depolarizza $\rightarrow$ le cariche positive entrate viaggiano per circa
    1--1,5\,mm (lunghezza dell'internodo) fino al nodo successivo, depolarizzandolo;
  \item nel frattempo le cariche positive esterne ripolarizzano il nodo depolarizzato per primo;
  \item il potenziale ``salta'' lunghi tratti di assone $\rightarrow$ conduzione molto più veloce,
    fino a 120\,m/s.
\end{itemize}

\subsection{La guaina mielinica}

\textbf{Guaina mielinica}: costituita da \textbf{cellule della glia} (cellule di supporto del
sistema nervoso) che si avvolgono attorno all'assone, circa un centinaio per ogni tratto,
ispessendo lo strato lipidico.
\begin{itemize}
  \item l'assone risulta elettricamente \textbf{isolato} dall'ambiente extracellulare nei tratti
    rivestiti;
  \item la guaina si interrompe solo a livello dei nodi di Ranvier ($\phi$ del nodo $\sim$3\,$\mu$m,
    internodo $\sim$1,5\,mm);
  \item i \textbf{canali ionici} per $Na^+$ e $K^+$ sono molto più abbondanti proprio a livello dei
    nodi di Ranvier: il neurone forma i canali dove servono e non ne dispone lungo i tratti
    mielinizzati, dove non è prevista conduzione.
\end{itemize}

\rubrica{Degenerazione della guaina}
Malattie demielinizzanti (p.\,es.\ \textbf{sclerosi multipla}) distruggono la guaina in alcuni punti
dell'assone, danneggiando la propagazione dei segnali elettrici e quindi delle informazioni.
\begin{itemize}
  \item le cariche non raggiungono più direttamente il nodo di Ranvier successivo: nei tratti
    demielinizzati la conduzione ridiventa punto a punto $\rightarrow$ \textbf{rallentamento};
  \item aggravante: dove era prevista la guaina i canali ionici sono molto scarsi, il che rallenta
    ulteriormente il passaggio dell'informazione.
\end{itemize}

% ---------------------------------------------------------------------------
\section{La sinapsi}

\textbf{Trasmissione dell'informazione}: il potenziale d'azione, che contiene l'informazione, deve
essere trasmesso ad altre cellule.
\begin{itemize}
  \item il potenziale si forma a livello del corpo cellulare e viaggia lungo l'assone, nodo per
    nodo, fino alla terminazione assonica, che perde la guaina mielinica;
  \item la trasmissione avviene attraverso il rilascio di molecole dette
    \textbf{neurotrasmettitori};
  \item \textbf{cellule bersaglio}: altre cellule del sistema nervoso, cellule muscolari e --- nel
    caso del sistema nervoso autonomo --- ghiandole o muscolatura liscia (p.\,es.\ del sistema
    gastroenterico).
\end{itemize}
